\chapter{Cadrer le besoin}\label{ch:cadrage}

\begin{objectifs}[Objectifs --- étape 0, \texttt{forge-cadrage}]
\begin{itemize}
  \item formuler une vision produit et des objectifs \emph{mesurables} ;
  \item construire une impact map reliant pourquoi, qui, comment et quoi ;
  \item prioriser avec MoSCoW et écarter ce qui ne sert aucun objectif ;
  \item reconnaître un chiffre inconnu et le noter \q{À mesurer} plutôt que de l'inventer.
\end{itemize}
\end{objectifs}

\begin{situation}[Sur le terrain : \sika{}]
Retour chez la présidente. Cette fois, avant tout écran, on pose des questions : combien de
groupes ? combien de membres ? combien de cotisations sont contestées chaque semaine ? La
trésorière générale a compté sur trois mois : 12\,\% des cotisations ont été discutées. Elle ne sait en
revanche pas combien de temps s'écoule entre la dernière cotisation et la remise du pot. Deux
objectifs émergent --- l'un avec un chiffre, l'autre sans.
\end{situation}

\section{Le concept : relier chaque fonctionnalité à un objectif}
\begin{concept}[Objectif mesurable]
Un objectif qui contient un \emph{indicateur}, une \emph{valeur actuelle}, une \emph{cible} et une
\emph{échéance}. \q{Améliorer la confiance} n'est pas mesurable ; \q{passer de 12\,\% à moins de 2\,\% de
cotisations contestées en six mois} l'est.
\end{concept}

\begin{concept}[Impact map]
Une carte à quatre niveaux qui relie le \emph{pourquoi} (l'objectif) au \emph{quoi} (la fonctionnalité)
par le \emph{qui} (l'acteur) et le \emph{comment} (le changement de comportement visé). Sa règle
d'or : une fonctionnalité qui ne remonte à aucun objectif sort du périmètre.
\end{concept}

\section{Pas à pas avec \texttt{forge-cadrage}}
Le skill vous interroge \emph{une question à la fois}, dans cet ordre :
\begin{enumerate}
  \item pour qui, et quel problème ;
  \item la situation actuelle, le cadre réglementaire (APDP, OHADA, marchés publics), les contraintes terrain (connectivité, terminaux, langues) ;
  \item les objectifs mesurables : indicateur, valeur actuelle, cible, échéance ;
  \item les acteurs, puis les changements de comportement ;
  \item enfin les fonctionnalités.
\end{enumerate}
Il renseigne ensuite \cmd{vision.md}, puis l'impact map (un diagramme et un tableau de priorisation).
Le livrable passe \emph{En revue} ; la validation revient au \textbf{sponsor client}, pas à l'ingénieur.

\section{Ce que \sika{} produit}
\subsection*{La vision}
\begin{code}
\begin{Verbatim}
Pour les groupes de tontine de l'Association Entraide
Dantokpa qui rapprochent chaque semaine un cahier et des
espèces, Sika est une plateforme de cotisation par Mobile
Money qui rend chaque cotisation vérifiable. Contrairement
au cahier, chaque cotisation y est horodatée et confirmée
par l'opérateur.
\end{Verbatim}
\end{code}

\subsection*{Les objectifs}
\noindent\begin{tabularx}{\linewidth}{lXlll}
\toprule
\textbf{\#} & \textbf{Objectif et indicateur} & \textbf{Actuel} & \textbf{Cible} & \textbf{Échéance} \\ \midrule
O1 & Rendre les cotisations indiscutables : \% de cotisations contestées & 12\,\% & $<2\,\%$ & 6 mois \\
O2 & Remettre le pot plus vite : délai de remise du pot & \textit{À mesurer} & $\le 24$\,h & 6 mois \\
\bottomrule
\end{tabularx}

\begin{piege}[Piège : inventer le chiffre qui manque]
Pour O2, personne ne connaît la valeur actuelle. Écrire \q{3 jours} \emph{parce que ça semble juste}
fausse toute la mesure future. Le skill note \emph{À mesurer} ; la première tâche du projet sera de
mesurer.
\end{piege}

\subsection*{L'impact map}
\begin{figure}[H]
\centering
\resizebox{\linewidth}{!}{%
\begin{tikzpicture}[
  n/.style={draw,rounded corners=3pt,align=center,font=\sffamily\small,minimum height=11mm,text width=3.2cm,inner sep=3pt},
  g/.style={n,fill=accent,draw=accent,text=white},
  a/.style={n,fill=accentl!20,draw=accentl},
  i/.style={n,fill=ochre!12,draw=ochre},
  d/.style={n,fill=leaf!12,draw=leaf},
  ar/.style={-{Stealth},thick,grey},
  h/.style={font=\sffamily\bfseries\scriptsize,text=grey}]
  \node[h] at (0,2.3) {POURQUOI}; \node[h] at (4.3,2.3) {QUI}; \node[h] at (8.6,2.3) {COMMENT}; \node[h] at (12.9,2.3) {QUOI};
  \node[g] (o1) at (0,0.9) {O1\\cotisations contestées $<2\,\%$};
  \node[g] (o2) at (0,-1.9) {O2\\pot remis en $\le 24$\,h};
  \node[a] (m) at (4.3,1.3) {Membre};
  \node[a] (t) at (4.3,-1.3) {Trésorière de groupe};
  \node[i] (i1) at (8.6,1.3) {Cotise depuis son téléphone, avec preuve};
  \node[i] (i2) at (8.6,-1.3) {Suit les cotisations sans pointer le cahier};
  \node[d] (d1) at (12.9,1.7) {Cotiser par Mobile Money};
  \node[d] (d2) at (12.9,0.4) {Solde du pot en temps réel};
  \node[d] (d3) at (12.9,-1.9) {Remettre le pot au bénéficiaire};
  \draw[ar] (o1)--(m); \draw[ar] (o1)--(t); \draw[ar] (o2)--(t);
  \draw[ar] (m)--(i1); \draw[ar] (t)--(i2);
  \draw[ar] (i1)--(d1); \draw[ar] (i1)--(d2); \draw[ar] (i2)--(d3);
\end{tikzpicture}}
\caption{L'impact map de \sika{} : chaque fonctionnalité remonte à un objectif.}
\end{figure}

\subsection*{La priorisation MoSCoW}
\begin{concept}[MoSCoW]
Quatre niveaux de priorité : \textbf{M}ust (indispensable au lot), \textbf{S}hould (important, pas bloquant),
\textbf{C}ould (utile si le temps le permet), \textbf{W}on't (explicitement hors lot).
\end{concept}
\noindent\begin{tabularx}{\linewidth}{Xcll}
\toprule
\textbf{Fonctionnalité} & \textbf{Objectif} & \textbf{Priorité} & \textbf{Lot} \\ \midrule
Cotiser par Mobile Money & O1 & Must & Lot 1 \\
Solde du pot en temps réel & O1 & Should & Lot 1 \\
Remettre le pot au bénéficiaire & O2 & Must & Lot 2 \\
Rappel SMS avant l'échéance & O1 & Could & Lot 2 \\
Messagerie vocale intégrée & \textit{aucun} & Won't & hors périmètre \\
\bottomrule
\end{tabularx}
La messagerie vocale, séduisante, ne sert ni O1 ni O2 : le skill la signale \emph{hors périmètre}. C'est exactement
ce que cette étape est faite pour détecter.

\begin{vousjouez}[À vous de jouer]
Sur votre projet : écrivez deux objectifs au format indicateur / actuel / cible / échéance. Pour chaque
chiffre que vous ne connaissez pas, écrivez \emph{À mesurer} et nommez qui peut le mesurer. Dessinez
ensuite une impact map de cinq nœuds. Y a-t-il une fonctionnalité \q{évidente} qui n'en fait pas partie ?
\end{vousjouez}

\begin{retenir}[À retenir]
\begin{itemize}
  \item Un objectif = indicateur + valeur actuelle + cible + échéance ; sinon il n'est pas mesurable.
  \item Pas d'objectif \texttt{O\#}, pas de fonctionnalité.
  \item Inconnu $\Rightarrow$ \emph{À mesurer}, jamais un chiffre inventé.
  \item Le sponsor client valide le cadrage.
\end{itemize}
\end{retenir}
\source{skills/forge-cadrage/SKILL.md}

\begin{acquis}
  \item Pourquoi \q{améliorer l'expérience des membres} n'est-il pas un objectif exploitable ?
  \item Que fait le skill d'une fonctionnalité sans objectif \texttt{O\#} ?
  \item Qui valide le cadrage ?
\end{acquis}
